% ^^A documentation part start with %
% ^^A All other parts are called definition parts
% ^^A Comments in documentation part
% ^^A equivalently use \iffalse ... \fi
% \iffalse meta-comment
% !TeX program = XeLaTeX
% !TeX encoding = UTF-8
%
% Copyright (C) 2016 by Van Abel <van141.abel(AT)gmail.com>
% ---------------------------------------------------------------------
%
% This work may be distributed and/or modified under the
% conditions of the LaTeX Project Public License, either
% version 1.3c of this license or (at your option) any later
% version. This version of this license is in
%    http://www.latex-project.org/lppl/lppl-1-3c.txt
% and the latest version of this license is in
%    http://www.latex-project.org/lppl.txt
% and version 1.3 or later is part of all distributions of
% LaTeX version 2005/12/01 or later.
%
% This work has the LPPL maintenance status `maintained'.
%
% The Current Maintainers of this work is Van Abel.
%
% ---------------------------------------------------------------------
%
%<*internal>
\iffalse
%</internal>
%<*readme>
---------------------------------------------------------------
mnsfc --- A package for NSFC Proposal
===============================================================
Released under the LaTeX Project Public License v1.3c or later
See http://www.latex-project.org/lppl.txt
---------------------------------------------------------------
This package is developed when I prepare the `NSFC` Proposal.

USAGE
---------------------------------------------------------------
Download the zip file and unzip to a local directory, then have a look at the
`mnsfc-main.pdf`, you can edit and compile it by 
`latexmk -xelatex -shell-escape mnsfc-main.tex`.
%</readme>
%<*internal>
\fi
\def\nameoflatex{plain}
\ifx\nameoflatex\fmtname\else
  \expandafter\begingroup
\fi
%</internal>
%
%<*install>

\input docstrip.tex
\keepsilent
\askforoverwritefalse
\preamble
----------------------------------------------------------------
mnsfc --- 数学类国家自然科学基金写作宏包
===============================================================
E-mail: van141.abel(AT)gmail.com
Released under the LaTeX Project Public License v1.3c or later
See http://www.latex-project.org/lppl.txt
----------------------------------------------------------------

\endpreamble
\postamble

Copyright (C) 2009 by You <van141.abel(AT)gmail.com>

This work may be distributed and/or modified under the
conditions of the LaTeX Project Public License (LPPL), either
version 1.3c of this license or (at your option) any later
version.  The latest version of this license is in the file:

http://www.latex-project.org/lppl.txt

This work is "maintained" (as per LPPL maintenance status) by
You.

This work consists of the file  mnsfc.dtx
and the derived files           mnsfc.ins,
                                mnsfc-main.tex,
                                mnsfc-main.pdf,
                                mnsfc.pdf and
                                mnsfc.sty.

\endpostamble
\generate{
  \usedir{tex/latex/mnsfc}
  \file{\jobname.sty}{\from{\jobname.dtx}{package}}
  \nopreamble\nopostamble
  \file{\jobname-main.tex}{\from{\jobname.dtx}{example-main}}
  \file{\jobname-refs.bib}{\from{\jobname.dtx}{example-bib}}
  \usedir{doc/latex/mnsfc}
  \file{README.txt}{\from{\jobname.dtx}{readme}}
}
\obeyspaces
\Msg{****************************************************}
\Msg{*                                                   }
\Msg{* To finish the installation you have to move the   }
\Msg{* following file into a directory searched by TeX   }
\Msg{* e.g. tex/latex/mnsfc:                          }
\Msg{*                                                   }
\Msg{* \jobname.sty                                      }
\Msg{*                                                   }
\Msg{* To produce the documentation run the file         }
\Msg{* \jobname.dtx through XeLaTeX.                     }
\Msg{*                                                   }
\Msg{* To produce the sample file run the file           }
\Msg{* \jobname-main.tex                                 }
\Msg{* through XeLaTeX and BibTeX                        }
\Msg{*                                                   }
\Msg{* Happy TeXing!                                     }
\Msg{*                                                   }
\Msg{****************************************************}
%</install>
%<install>\endbatchfile
%
%<*internal>
\usedir{source/latex/mnsfc}
\generate{
  \file{\jobname.ins}{\from{\jobname.dtx}{install}}
}
% ^^A No extra text add by DocStrip
\nopreamble\nopostamble
\usedir{doc/latex/mnsfc}
\generate{
  \file{README.txt}{\from{\jobname.dtx}{readme}}
}
% ^^A if xetex then end process of DocStrip by \endbatchfile
\ifx\nameoflatex\fmtname
  \expandafter\endbatchfile
\else
  % ^^A for xelatex we close the group
  \expandafter\endgroup
\fi
%</internal>
%
%<*driver>
\ProvidesFile{\jobname.dtx}
%</driver>
%<package>\NeedsTeXFormat{LaTeX2e}[2005/12/01]
%<package>\ProvidesPackage{mnsfc}
%<*package>
[2026/01/06 v.1.0.0 一个数学类国家自然科学基金写作宏包]
%</package>
%
%
%<*driver>
\documentclass{ltxdoc} %[draft]
% 需要使用 xelatex 编译以支持中文字符
\usepackage{palatino}
\usepackage[zihao=-4]{ctex}
\usepackage[margin=1in, left=1.2in, includeheadfoot]{geometry}
\usepackage[numbered]{hypdoc}
\usepackage{metalogo}
\EnableCrossrefs
\renewcommand\indexname{索引}
\IndexPrologue{%
  \section*{\indexname}
  \textit{
    斜体数字表示相应条目描述的页码, 
    而下划线的数字表示表示相应条目定义的页码 
    使用条目的页码用罗马数字表示. 
  }
}
\CodelineIndex
\RecordChanges
\def\glossaryname{版本历史}
\GlossaryPrologue{\section*{\glossaryname}}
\begin{document}
  \DocInput{\jobname.dtx}
  \PrintChanges
\PrintIndex
\end{document}
%</driver>
% \fi
%
% \CharacterTable
%  {Upper-case    \A\B\C\D\E\F\G\H\I\J\K\L\M\N\O\P\Q\R\S\T\U\V\W\X\Y\Z
%   Lower-case    \a\b\c\d\e\f\g\h\i\j\k\l\m\n\o\p\q\r\s\t\u\v\w\x\y\z
%   Digits        \0\1\2\3\4\5\6\7\8\9
%   Exclamation   \!     Double quote  \"     Hash (number) \#
%   Dollar        \$     Percent       \%     Ampersand     \&
%   Acute accent  \'     Left paren    \(     Right paren   \)
%   Asterisk      \*     Plus          \+     Comma         \,
%   Minus         \-     Point         \.     Solidus       \/
%   Colon         \:     Semicolon     \;     Less than     \<
%   Equals        \=     Greater than  \>     Question mark \?
%   Commercial at \@     Left bracket \[     Backslash     \\
%   Right bracket \]     Circumflex   \^     Underscore    \_
%   Grave accent  \`     Left brace    \{     Vertical bar  \|
%   Right brace   \}     Tilde         \~}
%
%\GetFileInfo{\jobname.dtx}
%
%\DoNotIndex{\%,\#,\$,\%,\&,\\,\^,\_,\~,\[,\],\',\`,\documentclass}
%\DoNotIndex{\begin,\end,\par,\def,\let,\gdef,\expandafter,\ifx,\ifnum,\else,\fi,\relax,\empty}
%\DoNotIndex{\newcommand,\RequirePackage,\usepackage,\renewcommand,\renewenvironment}
%\DoNotIndex{\section,\subsection,\subsubsection,\paragraph,\item,\label,\cite,\cites,\eqref,\href}
%\DoNotIndex{\emph,\textbf,\textcolor,\textup,\noindent,\vspace,\hspace,\clearpage,\newpage,\vfill}
%\DoNotIndex{\zhlipsum,\bibliography,\bibliographystyle,\nocite}
%\DoNotIndex{\@ifpackageloaded,\@ifundefined,\@nil,\@oldsection,\@oldsubsection}
%\DoNotIndex{\@plus,\@pnumwidth,\@secpenalty,\@tempdima}
%\DoNotIndex{\ctexset,\zihao,\songti,\kaishu,\chinese,\ccwd}
%\DoNotIndex{\geometry,\hypersetup,\setlist,\setlength,\setcounter,\newcounter,\newif}
%\DoNotIndex{\RenewDocumentCommand,\NewEnviron,\AfterEndEnvironment,\AtBeginDocument}
%\DoNotIndex{\IfFileExists,\IfNoValueTF,\PackageWarning,\definecolor,\setmainfont}
%\DoNotIndex{\pagestyle,\fancyhf,\headrulewidth,\footrulewidth}
%\DoNotIndex{\ifmnsfc@final,\ifmnsfc@debug}
%\DoNotIndex{\value,\number,\numexpr,\the,\arabic,\alph,\roman,\z@,\p@}
%\DoNotIndex{\baselineskip,\parindent,\parskip,\textwidth,\leftskip,\rightskip,\parfillskip}
%\DoNotIndex{\clubpenalty,\widowpenalty,\displaywidowpenalty}
%\DoNotIndex{\linenumbers,\nolinenumbers,\modulolinenumbers,\linenumberfont,\linenumbersep}
%\DoNotIndex{\c@linenumber,\c@tocdepth,\getpagerefnumber}
%\DoNotIndex{\addpenalty,\addvspace,\advance,\begingroup,\endgroup,\leavevmode,\nobreak}
%\DoNotIndex{\hfil,\hss,\hb@xt@,\hskip,\bfseries,\normalfont,\small,\Large,\LARGE}
%\DoNotIndex{\color,\red,\blue,\bullet,\circ,\space,\allowbreak,\typeout,\string}
%\DoNotIndex{\raggedright,\ldots,\int,\,,\(,\)}
%\DoNotIndex{\CTEX@section@tocline,\CTEX@subsection@tocline,\CTEX@subsubsection@tocline}
%\DoNotIndex{\CTEX@paragraph@tocline,\l@section,\l@subsection,\l@subsubsection,\l@paragraph}
%\DoNotIndex{\numberline}
%\DoNotIndex{\mnsfc@applicant,\mnsfc@check@colon,\mnsfc@check@colon@cn,\mnsfc@check@colon@en}
%\DoNotIndex{\mnsfc@debugtrue,\mnsfc@finalfalse,\mnsfc@finaltrue}
%\DoNotIndex{\mnsfc@kwen,\mnsfc@kwzh,\mnsfc@projecttitle,\mnsfc@researchtype,\mnsfc@toc@title,\mnsfc@unit}
%\DoNotIndex{\BODY,\DeclareOption,\ExecuteOptions,\ProcessOptions,\setstretch,\stepcounter}
%\expandafter\DoNotIndex\expandafter{\string\&}
%\expandafter\DoNotIndex\expandafter{\string\{}
%\expandafter\DoNotIndex\expandafter{\string\}}
%
%\def\DescribeOption{\leavevmode
%  \begingroup\MakePrivateLetters\Describe@Option}
%\def\Describe@Option#1{\endgroup%
%  \marginpar{\raggedleft\PrintDescribeOption{#1}}%
%  \SpecialOptionIndex{#1}\ignorespaces}
%\def\PrintDescribeOption#1{\strut\MacroFont #1}
%\def\SpecialOptionIndex#1{%
%  \index{\quotechar/#1% sort options as `Symbols'
%  \actualchar{\protect\ttfamily#1}\encapchar usage}%
%  \CatIndex{#1}{option}}
%\def\CatIndex#1#2{\index{#1\actualchar{\protect\ttfamily #1} (#2)\encapchar hyperpage}}
%
%\title{^^A
%   \textsf{mnsfc} --- 数学类国家自然科学基金写作宏包\thanks{^^A 
%     这是对版本号为\fileversion 的文档说明, 最后修改日期为 \filedate.^^A
%  }^^A
%}
%\author{^^A
%  Van Abel\thanks{E-mail: van141.abel(AT)gmail.com}^^A
%}
%\date{发布日期 \filedate}
%
%\maketitle
% \begin{abstract}
% 数学类国家自然科学基金申请书 LaTeX 模板，支持 draft 和 final 双模式，提供行号、字数统计等功能。
% \end{abstract}
%\changes{v1.0.0}{2026/01/06}{初始版本}
% \tableofcontents
%
% \section{简介}
%
% \textsf{mnsfc} 是一个用于编写数学国家自然科学基金申请书的 LaTeX 模板包。它提供了符合 NSFC 规范的页面布局、字体设置和格式要求，同时支持 draft 和 final 两种模式，方便用户在写作和提交之间切换。
%
% \section{简明使用教程}
%
% \subsection{基本设置}
%
% 在文档中加载包：
% \begin{verbatim}
%
% \documentclass[UTF8,zihao=-4]{ctexart} %for macos please set: fontset=macnew
% \usepackage[final]{mnsfc} %draft is default, final is for the final version
% \end{verbatim}
% 注意： 苹果系统需要设置|fontset=macnew|来正常显示楷体加粗。
%\section{选项、命令以及环境}%
%\subsection{模式切换}
% 选项可以通过传递给文档类或者宏包的形式启用 %
%
%\DescribeOption{draft}
% |draft| 默认选项：草稿版，输出标题、摘要以及字数统计等信息。
%
%\DescribeOption{final}
% |final| 最终提交版，不输出标题、摘要以及字数统计等信息。
%
%\DescribeOption{debug}
% |debug| 输出debug信息，一般只针对开发者使用，普通用户不建议使用。
%
% \subsection{项目基本信息设置}
%
% \begin{verbatim}
% \mnsfcProjectTitle{（填写项目名称）}
% \mnsfcKeywordsZH{关键词1；关键词2；关键词3}
% \mnsfcKeywordsEN{keywords1; keywords2; keywords3}
% % 选择研究属性（二选一）：
% % \mnsfcResearchTypeFree  % 自由探索类基础研究
% % \mnsfcResearchTypeTarget % 目标导向类基础研究
% % 或者使用自定义文本：
% % \mnsfcResearchType{自由探索类基础研究}
% % \mnsfcResearchType{目标导向类基础研究}
% \end{verbatim}
%
%\section{功能特性}
%
%\subsection{双模式支持}
%
%\DescribeOption{draft}
%|draft| 模式（默认）：
%\begin{itemize}
%  \item 显示标题页（项目名称）
%  \item 显示中文和英文摘要
%  \item 显示写作提示（|\mnsfcNote|）
%  \item 显示行号（每 5 行显示一次）
%  \item 显示字数统计页面
%  \item 显示 |zhlipsum| 占位文本
%\end{itemize}
%
%\DescribeOption{final}
%|final| 模式：
%\begin{itemize}
%  \item 隐藏标题页
%  \item 隐藏摘要内容
%  \item 隐藏所有写作提示
%  \item 隐藏行号
%  \item 隐藏字数统计
%  \item 隐藏占位文本
%  \item 仅显示符合 NSFC 规范的正文内容
%\end{itemize}
%
%\DescribeOption{debug}
%|debug| 选项：输出调试信息，一般只针对开发者使用，普通用户不建议使用。
%
%\subsection{页面布局与格式}
%
%\textbf{页面设置}：
%\begin{itemize}
%  \item A4 纸张
%  \item 页边距：左右 30mm，上下 25mm
%  \item 无页眉页脚
%\end{itemize}
%
%\textbf{正文字体与格式}：
%\begin{itemize}
%  \item 字体：宋体（|\songti|）
%  \item 字号：小四号（|\zihao{-4}|）
%  \item 行距：22 磅（|\baselineskip=22pt|）
%  \item 首行缩进：0.74 cm
%  \item 段落间距：0.45em
%\end{itemize}
%
%\textbf{分页控制}：
%\begin{itemize}
%  \item 避免孤行和寡行（|\clubpenalty=300|, |\widowpenalty=300|）
%\end{itemize}
%
%\subsection{标题格式}
%
%\DescribeMacro{\section}
%\cs{section} 命令：
%\begin{itemize}
%  \item 格式：|（一）|、|（二）|...（中文数字）
%  \item 字体：楷体四号，蓝色，加粗
%  \item 支持可选参数：|\section[说明文字]{标题}|
%  \item 自动跳过行号
%\end{itemize}
%
%\DescribeMacro{\subsection}
%\cs{subsection} 命令：
%\begin{itemize}
%  \item 格式：|1.|、|2.|...（阿拉伯数字）
%  \item 字体：楷体四号，蓝色，加粗
%  \item 支持可选参数：|\subsection[说明文字]{标题}|
%  \item 自动跳过行号
%\end{itemize}
%
%\DescribeMacro{\subsubsection}
%\cs{subsubsection} 命令：
%\begin{itemize}
%  \item 格式：|(1)|、|(2)|...（带括号的阿拉伯数字）
%  \item 字体：宋体小四号，加粗
%  \item 自动跳过行号
%\end{itemize}
%
%\DescribeMacro{\paragraph}
%\cs{paragraph} 命令：
%\begin{itemize}
%  \item 格式：不带编号
%  \item 字体：宋体小四号，加粗
%  \item 自动跳过行号
%\end{itemize}
%
%\subsection{列表环境}
%
%\DescribeEnv{enumerate}
%\cs{enumerate} 环境（有序列表）：
%\begin{itemize}
%  \item 第一层：|1)|, |2)|, |3)|...
%  \item 第二层：|a)|, |b)|, |c)|...
%  \item 第三层：|i)|, |ii)|, |iii)|...
%\end{itemize}
%
%\DescribeEnv{itemize}
%\cs{itemize} 环境（无序列表）：
%\begin{itemize}
%  \item 第一层：实心圆点 |$\bullet$|
%  \item 第二层：空心圆点 |$\circ$|
%  \item 第三层：短横线 |-|
%\end{itemize}
%
%所有列表项使用宋体小四号，首行缩进 0.74cm。
%
%\subsection{行号功能（仅 \cs{draft} 模式）}
%
%\begin{itemize}
%  \item 自动在正文开始后启用行号
%  \item 每 5 行显示一次行号（|mod=5|）
%  \item 自动跳过所有标题行（\cs{section}/\cs{subsection}/\cs{subsubsection}/\cs{paragraph}）
%  \item 行号字体：小四号无衬线字体
%  \item 使用 \cs{mnsfcExcludeLine}\marg{内容} 可以排除短行（如 "无。"）的行号显示
%\end{itemize}
%
%\subsection{字数统计（仅 \cs{draft} 模式）}
%
%\textbf{基于行号计算}（如果启用行号）：
%\begin{itemize}
%  \item 计算公式：|字数 = (行号 - 1) × 35字/行|
%  \item 每行按 35 个汉字（包括符号）计算
%  \item 自动排除使用 \cs{mnsfcExcludeLine} 标记的短行
%\end{itemize}
%
%\textbf{基于页数估算}（如果未启用行号）：
%\begin{itemize}
%  \item 计算公式：|字数 ≈ 页数 × 35行/页 × 35字/行|
%\end{itemize}
%
%显示信息：
%\begin{itemize}
%  \item 正文总页数（不含封面和摘要）
%  \item 正文行数
%  \item 已排除短行数（如果有）
%  \item 正文字数估算
%\end{itemize}
%
%字数统计页面不显示行号。
%
%\subsection{参考文献支持}
%
%\begin{itemize}
%  \item 使用 |amsrefs| 包管理参考文献
%  \item 提供 \cs{mnsfcReferences}\marg{文件名} 命令
%  \item 自动居中显示"参考文献"标题
%  \item 使用楷体四号标题，宋体小四号正文
%  \item 支持 MR 链接（需要参考文献数据中包含 |MRNUMBER| 字段）
%\end{itemize}
%
%\subsection{其他功能}
%
%\textbf{项目信息}：
%\begin{itemize}
%  \item \cs{mnsfcProjectTitle}\marg{项目名称}
%  \item \cs{mnsfcKeywordsZH}\marg{关键词1；关键词2；关键词3}
%  \item \cs{mnsfcKeywordsEN}\marg{keywords1; keywords2; keywords3}
%  \item \cs{mnsfcResearchTypeFree} 或 \cs{mnsfcResearchTypeTarget}（研究属性）
%\end{itemize}
%
%\textbf{写作提示}：\cs{mnsfcNote}\marg{提示内容}：仅在 |draft| 模式下显示灰色提示文字
%
%\textbf{标题页}：\cs{mnsfcTitlePage}：仅在 |draft| 模式下显示
%
%\textbf{摘要环境}：
%\begin{itemize}
%  \item \begin{verbatim}\begin{mnsfcAbstractZH}...\end{mnsfcAbstractZH}：中文摘要\end{verbatim}
%  \item \begin{verbatim}\begin{mnsfcAbstractEN}...\end{mnsfcAbstractEN}：英文摘要\end{verbatim}
%  \item 仅在 |draft| 模式下显示，|final| 模式下自动隐藏
%\end{itemize}
%
%\textbf{自动插入正文标题}：
%\begin{itemize}
%  \item 在英文摘要结束后自动插入"报告正文（2026版）"标题
%  \item 自动插入格式说明文字
%\end{itemize}
%
% \StopEventually{}
%
% \section{代码实现}
%
%    \begin{macrocode}
%<*package>
%% ==============================
%% mnsfc.sty  (final / draft)
%% ==============================
%% ---------- Options ----------
\newif\ifmnsfc@final
\newif\ifmnsfc@debug
\DeclareOption{final}{\mnsfc@finaltrue}
\DeclareOption{draft}{\mnsfc@finalfalse}
\DeclareOption{debug}{\mnsfc@debugtrue}
% 使用 \string 避免选项名称被索引
% 但这样会影响代码可读性，所以改用其他方法
\ExecuteOptions{draft} %% default
\ProcessOptions\relax

%% ---------- Packages ----------
\RequirePackage{geometry}
\RequirePackage{setspace}
\RequirePackage{fancyhdr}
\RequirePackage{lastpage}
\RequirePackage{refcount}
\RequirePackage{graphicx}
\RequirePackage{xcolor}
\RequirePackage{hyperref}
\RequirePackage{bookmark}
\RequirePackage{enumitem}
\RequirePackage{amsmath,amssymb}
\RequirePackage{microtype}
\RequirePackage{environ}
\RequirePackage{etoolbox}
%% amsrefs 不在这里加载，让用户在文档中加载以控制选项
%% \RequirePackage{amsrefs}
\RequirePackage{xparse}
\RequirePackage{totcount}
\RequirePackage{lineno}


%% ---------- Page layout (按 NSFC 正文习惯) ----------
\geometry{
  a4paper,
  left=30mm,
  right=30mm,
  top=25mm,
  bottom=25mm
}
%\setstretch{1.25}
%% 定义 NSFC 常用颜色
\definecolor{NSFCRed}{RGB}{255,0,0}
\definecolor{NSFCBlue}{RGB}{0,112,192}
\newcommand{\red}[1]{\textcolor{NSFCRed}{#1}}
\newcommand{\blue}[1]{\textcolor{NSFCBlue}{#1}}

%% 定义优雅的链接颜色（与 NSFC 蓝色协调）
\definecolor{NSFCLinkBlue}{RGB}{0,80,160}      %% 深蓝色，用于内部链接
\definecolor{NSFCCiteBlue}{RGB}{50,100,150}   %% 中等蓝色，用于引用
\definecolor{NSFCURLBlue}{RGB}{0,112,192}    %% NSFC 标准蓝色，用于 URL

%%% ---------- Times New Roman 字体设置 ----------
%% 设置英文字体为 Times New Roman
%% 使用 XeLaTeX 时，fontspec 已由 ctex 加载
\@ifpackageloaded{fontspec}{%
  \setmainfont{Times New Roman}[
    BoldFont = Times New Roman Bold,
    ItalicFont = Times New Roman Italic,
    BoldItalicFont = Times New Roman Bold Italic
  ]%
}{%
  %% 如果 fontspec 未加载（使用 pdfLaTeX），则使用 times 包
  \IfFileExists{times.sty}{%
    \RequirePackage{times}%
  }{%
    \PackageWarning{mnsfc}{Times New Roman font not available.%
      Using default font.}%
  }%
}

%% ---------- Hyperref ----------
\hypersetup{
  colorlinks=true,
  linkcolor=NSFCLinkBlue,      %% 深蓝色，优雅且专业
  citecolor=NSFCCiteBlue,      %% 中等蓝色，与链接区分但协调
  urlcolor=NSFCURLBlue,         %% NSFC 标准蓝色，用于 URL
  pdfborder={0 0 0}
}

%% ---------- Header / Footer ----------
\pagestyle{empty}   %% 使用 plain 样式
\fancyhf{}          %% 清空所有页眉页脚内容
\renewcommand{\headrulewidth}{0pt} %% 去掉页眉横线
\renewcommand{\footrulewidth}{0pt} %% 去掉页脚横线


%% ==============================
%% 正文字体与行距（NSFC final 规范）
%% ==============================

%\RequirePackage{ctex}

%% 小四号，宋体
%% 检查 ctex 包是否加载，如果未加载则给出警告
\@ifpackageloaded{ctex}{%
  \AtBeginDocument{%
    \zihao{-4}%
    \songti
    \setcounter{tocdepth}{1}%% 设置目录深度为 1（只显示 section）
  }%
}{%
  %% 如果 ctex 未加载，尝试延迟执行
  \AtBeginDocument{%
    \@ifundefined{zihao}{%
      \PackageWarning{mnsfc}{ctex package not loaded. Font settings (\zihao,%
        \songti) may not work. Please use \string\documentclass{ctexart}%
        or load ctex package.}%
    }{%
      \zihao{-4}%
    }%
    \@ifundefined{songti}{%
      \PackageWarning{mnsfc}{songti command not found. Please use %
        \string\documentclass{ctexart} or load ctex package.}%
    }{%
      \songti
    }%
    \setcounter{tocdepth}{1}%% 设置目录深度为 1（只显示 section）
  }%
}

%% 行距：22 磅（这是核心）
\setlength{\baselineskip}{22pt}

%% 首行缩进 0.74 cm
\setlength{\parindent}{0.74cm}
%% 段前 0 行，段后 0.5 行（= 11pt）
\setlength{\parskip}{0.45em}

%% 允许正常分页，避免孤行和寡行
\clubpenalty=300
\widowpenalty=300
\displaywidowpenalty=50

%% ---------- Metadata ----------
\newcommand{\mnsfcProjectTitle}[1]{\def\mnsfc@projecttitle{#1}}
\newcommand{\mnsfc@projecttitle}{（项目名称）}

\newcommand{\mnsfcApplicant}[1]{\def\mnsfc@applicant{#1}}
\newcommand{\mnsfcUnit}[1]{\def\mnsfc@unit{#1}}

\newcommand{\mnsfcKeywordsZH}[1]{\def\mnsfc@kwzh{#1}}
\newcommand{\mnsfcKeywordsEN}[1]{\def\mnsfc@kwen{#1}}

%% 研究属性：自由探索类基础研究 或 目标导向类基础研究
\newcommand{\mnsfcResearchType}[1]{%
  \def\mnsfc@researchtype{#1}%
}
\newcommand{\mnsfc@researchtype}{}%% 默认为空

%% 便捷命令
\newcommand{\mnsfcResearchTypeFree}{\mnsfcResearchType{自由探索类基础研究}}
\newcommand{\mnsfcResearchTypeTarget}{\mnsfcResearchType{目标导向类基础研究}}

%% ---------- Draft-only note ----------
\newcommand{\mnsfcNote}[1]{%
  \ifmnsfc@final\relax%
  \else{\par\textcolor{gray}{\small【写作提示】#1}\par}%
  \fi%
}
%% ---------- Title ----------
\newcommand{\mnsfcTitlePage}{%
\begin{center}
  \ifmnsfc@final
  \else
  {\Large\bfseries 数学类国家自然科学基金面上项目申请书\par}
  \vspace{1em}
  {\LARGE\bfseries \mnsfc@projecttitle\par}
  \vspace{0.5em}
  \ifx\mnsfc@researchtype\empty
    %% 如果未设置研究属性，显示提示（左对齐，与中文摘要一致）
    \begin{minipage}{\textwidth}\raggedright
      {\normalfont\kaishu\zihao{4}\textcolor{NSFCRed}{（请选择研究属性：自由探索类基础研究 / 目标导向类基础研究）}\par}%
    \end{minipage}\par
  \else
    %% 显示研究属性（左对齐，与中文摘要一致）
    \begin{minipage}{\textwidth}\raggedright
      {\normalfont\textbf{研究属性：}\mnsfc@researchtype}%
    \end{minipage}\par
  \fi
  \vspace{0.5em}
  \fi
\end{center}
\vspace{-\baselineskip}%% 减少标题页后的间距
}

%% ---------- Abstracts (final only) ----------

\ifmnsfc@final

  %% ========== FINAL：整块吞掉 ==========
  \NewEnviron{mnsfcAbstractZH}{}
  \NewEnviron{mnsfcAbstractEN}{}

\else

  %% ========== DRAFT：正常摘要 ==========
  \NewEnviron{mnsfcAbstractZH}{%
    \par\noindent\normalfont\textbf{中文摘要：}\small
    \BODY\par
    \noindent\normalfont\textbf{关键词：}\mnsfc@kwzh\par
  }
  \NewEnviron{mnsfcAbstractEN}{%
    \par\noindent\textbf{Abstract:}\small
    \BODY\par
    \noindent\textbf{Keywords: }\mnsfc@kwen\par
  }

\fi

%% 在英文摘要结束后自动插入 NSFCP
\AfterEndEnvironment{mnsfcAbstractEN}{%
\ifmnsfc@final\vspace{-1.7\baselineskip}\fi% 减少标题后的间距
\begin{center}
  {\kaishu \zihao{3} \textbf{报告正文（\( 2026 \)版）}}
\end{center}
\vspace{-0.7\baselineskip}% 减少标题后的间距
\kaishu\zihao{4}\hspace*{2\ccwd}参照以下提纲撰写，要求内容翔实、清晰，层次分明，标题突出。\red{\textbf{申请书正文原则上不超过\( 30 \)页，鼓励简洁表达。}}%
  \textbf{\blue{请勿删除或改动下述提纲标题及括号中的文字。}}
\par%
\songti\zihao{-4}% 恢复正文为宋体小四号
%% 在正文开始后启用行号（仅在 draft 模式下）
\ifmnsfc@final\relax%
\else%
  \modulolinenumbers[5]%% 每 5 行显示一次行号
  \linenumbers%
  %% 配置行号：跳过 section/subsection 等标题
  %%\setlength\linenumbersep{3pt}%% 行号与文本的间距
  \renewcommand\linenumberfont{\normalfont\zihao{-4}}%% 行号字体
\fi%
}
\ctexset{
punct=quanjiao,
  section = {
    name = {},
    number = {}, %%\chinese{section}
    numberformat = {（\chinese{section}）},
    format = {
      \ifmnsfc@final\relax%
      \else\nolinenumbers\fi% 跳过 section 标题的行号
      \hspace*{0.74cm}
      \kaishu
      \zihao{4}
      \color{NSFCBlue}
      \bfseries
    },
    aftername = {\space},
    beforeskip = 0\baselineskip,
    afterskip  = 0\baselineskip,
    aftertitle = {%
      \ifmnsfc@final\relax%
      \else\linenumbers\fi% 恢复行号
    },
  },
  subsection = {
    name = {},
    number = {}, %%\chinese{section}
    numberformat = {\normalfont\kaishu\arabic{subsection}.}, 
    format = {
      \ifmnsfc@final\relax%
      \else\nolinenumbers\fi% 跳过 subsection 标题的行号
      \hspace*{0.74cm}
      \kaishu
      \zihao{4}
      \color{NSFCBlue}
      \bfseries
    },
    aftername = {\space},
    beforeskip = 0\baselineskip,
    afterskip  = 0\baselineskip,
    aftertitle = {%
      \ifmnsfc@final\relax%
      \else\linenumbers\fi% 恢复行号
    },
  },
  subsubsection = {
    name = {},
    number = {},
    numberformat = {（\arabic{subsubsection}）}, % (1)
    format = {
      \ifmnsfc@final\relax%
      \else\nolinenumbers\fi% 跳过 subsubsection 标题的行号
      \hspace*{0.74cm}
      \songti
      \zihao{-4}
      \bfseries
    },
    aftername = {\space},
    beforeskip = 0.3\baselineskip,
    afterskip  = 0.2\baselineskip,
    aftertitle = {%
      \ifmnsfc@final\relax%
      \else\linenumbers\fi% 恢复行号
    },
  },
  paragraph = {
    name = {},
    number = {}, %% 不带编号
    numberformat = {},
    format = {
      \ifmnsfc@final\relax%
      \else\nolinenumbers\fi% 跳过 paragraph 标题的行号
      \hspace*{0.74cm}
      \songti
      \zihao{-4}
      \bfseries
    },
    aftername = {\space},
    beforeskip = 0.2\baselineskip,
    afterskip  = 0.1\baselineskip,
    aftertitle = {%
      \ifmnsfc@final\relax%
      \else\linenumbers\fi% 恢复行号
    },
  }
}

%% ==================================================
%% 配置列表环境（enumerate 和 itemize）
%% ==================================================
%% 设置 enumerate 和 itemize 的缩进和格式
\setlist{
  leftmargin=1.5em,        %% 左缩进
  itemindent=0.74cm,       %% 首行缩进（与正文一致）
  labelsep=0.5em,          %% 标签与文本的间距
  topsep=0.3\baselineskip, %% 列表前后的间距
  partopsep=0pt,           %% 段落开始时的额外间距
  parsep=0.1\baselineskip, %% 列表项之间的间距
  itemsep=0.1\baselineskip %% 列表项之间的间距
}

%% enumerate 环境：使用阿拉伯数字编号
\setlist[enumerate,1]{
  label=\arabic*),          %% 第一层：1), 2), 3)...
  font=\songti\zihao{-4}
}
\setlist[enumerate,2]{
  label=\alph*),            %% 第二层：a), b), c)...
  font=\songti\zihao{-4}
}
\setlist[enumerate,3]{
  label=\roman*),           %% 第三层：i), ii), iii)...
  font=\songti\zihao{-4}
}

%% itemize 环境：使用中文符号
\setlist[itemize,1]{
  label=$\bullet$,          %% 第一层：实心圆点
  font=\songti\zihao{-4}
}
\setlist[itemize,2]{
  label=$\circ$,            %% 第二层：空心圆点
  font=\songti\zihao{-4}
}
\setlist[itemize,3]{
  label=$-$                 %% 第三层：短横线
}

%% ==================================================
%% 重新定义 \section 和 \subsection 以支持说明文字
%% ==================================================
%% 保存原始的 section 和 subsection 命令
\let\@oldsection\section
\let\@oldsubsection\subsection

%% 辅助命令：从标题中移除末尾的冒号（用于目录，支持中文冒号：和英文冒号:）
%% 使用字符串匹配，将结果存储在 \mnsfc@toc@title 中
%% 先尝试匹配中文冒号
\def\mnsfc@check@colon@cn#1：#2\@nil{%
  \ifx\relax#2\relax
    %% 没有中文冒号，尝试英文冒号
    \expandafter\mnsfc@check@colon@en#1:\@nil
  \else
    %% 有中文冒号，#1 是冒号前的内容
    \gdef\mnsfc@toc@title{#1}%
  \fi
}
%% 尝试匹配英文冒号
\def\mnsfc@check@colon@en#1:#2\@nil{%
  \ifx\relax#2\relax
    %% 没有英文冒号，使用全部内容
    \gdef\mnsfc@toc@title{#1}%
  \else
    %% 有英文冒号，#1 是冒号前的内容
    \gdef\mnsfc@toc@title{#1}%
  \fi
}
%% 主入口
\def\mnsfc@check@colon#1\@nil{%
  \mnsfc@check@colon@cn#1：\@nil
}

%% 修改目录格式：移除 \numberline {} 和格式命令
%% 重定义 ctex 的 tocline 命令，使其不添加 \numberline{}
\AtBeginDocument{%
  \ifmnsfc@debug
    \typeout{^^J[MNSFC] 目录格式修复已激活：移除 \numberline{} 和格式命令^^J}%
  \fi
  %% 重定义 \CTEX@section@tocline，不添加 \numberline{}
  \@ifundefined{CTEX@section@tocline}{}{%
    \renewcommand{\CTEX@section@tocline}[2]{#2}%
    \ifmnsfc@debug
      \typeout{[MNSFC] 已重定义 CTEX@section@tocline}%
    \fi
  }%
  %% 重定义 \CTEX@subsection@tocline，不添加 \numberline{}
  \@ifundefined{CTEX@subsection@tocline}{}{%
    \renewcommand{\CTEX@subsection@tocline}[2]{#2}%
    \ifmnsfc@debug
      \typeout{[MNSFC] 已重定义 CTEX@subsection@tocline}%
    \fi
  }%
  %% 重定义 \CTEX@subsubsection@tocline，不添加 \numberline{}
  \@ifundefined{CTEX@subsubsection@tocline}{}{%
    \renewcommand{\CTEX@subsubsection@tocline}[2]{#2}%
    \ifmnsfc@debug
      \typeout{[MNSFC] 已重定义 CTEX@subsubsection@tocline}%
    \fi
  }%
  %% 重定义 \CTEX@paragraph@tocline，不添加 \numberline{}
  \@ifundefined{CTEX@paragraph@tocline}{}{%
    \renewcommand{\CTEX@paragraph@tocline}[2]{#2}%
    \ifmnsfc@debug
      \typeout{[MNSFC] 已重定义 CTEX@paragraph@tocline}%
    \fi
  }%
  %% 重新定义 \l@section 以移除 \numberline{}
  \renewcommand{\l@section}[2]{%
    \ifnum \c@tocdepth >\z@
      \addpenalty\@secpenalty
      \addvspace{1.0em \@plus\p@}%
      \setlength\@tempdima{1.5em}%
      \begingroup
        \parindent \z@ \rightskip \@pnumwidth
        \parfillskip -\@pnumwidth
        \leavevmode \bfseries
        \advance\leftskip\@tempdima
        \hskip -\leftskip
        %% #1 是标题（不含 \numberline{}，因为我们通过可选参数传递了纯标题）
        #1\nobreak\hfil \nobreak\hb@xt@\@pnumwidth{\hss #2}\par
      \endgroup
    \fi
  }%
  %% 重新定义 \l@subsection 以移除 \numberline{}
  \renewcommand{\l@subsection}[2]{%
    \ifnum \c@tocdepth >1\relax
      \addpenalty\@secpenalty
      \addvspace{0.5em \@plus\p@}%
      \setlength\@tempdima{2.5em}%
      \begingroup
        \parindent \z@ \rightskip \@pnumwidth
        \parfillskip -\@pnumwidth
        \leavevmode
        \advance\leftskip\@tempdima
        \hskip -\leftskip
        %% #1 是标题（不含 \numberline{} 和格式命令）
        #1\nobreak\hfil \nobreak\hb@xt@\@pnumwidth{\hss #2}\par
      \endgroup
    \fi
  }%
  %% 重新定义 \l@subsubsection 以移除 \numberline{}
  \renewcommand{\l@subsubsection}[2]{%
    \ifnum \c@tocdepth >2\relax
      \addpenalty\@secpenalty
      \addvspace{0.3em \@plus\p@}%
      \setlength\@tempdima{3.5em}%
      \begingroup
        \parindent \z@ \rightskip \@pnumwidth
        \parfillskip -\@pnumwidth
        \leavevmode
        \advance\leftskip\@tempdima
        \hskip -\leftskip
        %% #1 是标题（不含 \numberline{}）
        #1\nobreak\hfil \nobreak\hb@xt@\@pnumwidth{\hss #2}\par
      \endgroup
    \fi
  }%
  %% 重新定义 \l@paragraph 以移除 \numberline{}
  \renewcommand{\l@paragraph}[2]{%
    \ifnum \c@tocdepth >3\relax
      \addpenalty\@secpenalty
      \addvspace{0.2em \@plus\p@}%
      \setlength\@tempdima{4.5em}%
      \begingroup
        \parindent \z@ \rightskip \@pnumwidth
        \parfillskip -\@pnumwidth
        \leavevmode
        \advance\leftskip\@tempdima
        \hskip -\leftskip
        %% #1 是标题（不含 \numberline{}）
        #1\nobreak\hfil \nobreak\hb@xt@\@pnumwidth{\hss #2}\par
      \endgroup
    \fi
  }%
}

% 重新定义 \section，支持 \section[说明]{标题} 格式
\RenewDocumentCommand{\section}{o m}{%
  \ifmnsfc@debug
    \typeout{[MNSFC] \section 命令已执行：标题="#2"}%
  \fi
  \ifmnsfc@final\relax%
  \else\nolinenumbers\fi% 跳过 section 标题的行号
  %% 无论是否有可选参数，都要移除标题中的冒号用于目录
  \expandafter\mnsfc@check@colon#2\@nil
  \ifmnsfc@debug
    \typeout{[MNSFC] 移除冒号：原标题="#2"，目录标题="\mnsfc@toc@title"}%
  \fi
  \IfNoValueTF{#1}{%
    % 没有说明文字
    \@oldsection[\mnsfc@toc@title]{#2}%
  }{%
    % 有说明文字
    \@oldsection[\mnsfc@toc@title]{#2}%
    \kaishu\zihao{4}\blue{#1}\par%
  }%
  \ifmnsfc@final\relax%
  \else\linenumbers\fi%% 恢复行号
  \songti\zihao{-4}%% 确保后续正文为宋体小四号
  \allowbreak
}

% 重新定义 \subsection，支持 \subsection[说明]{标题} 格式
\RenewDocumentCommand{\subsection}{o m}{%
  \ifmnsfc@debug
    \typeout{[MNSFC] \subsection 命令已执行：标题="#2"}%
  \fi
  \ifmnsfc@final\relax%
  \else\nolinenumbers\fi% 跳过 subsection 标题的行号
  \IfNoValueTF{#1}{%
    %% 没有说明文字
    \@oldsubsection{#2}%
  }{%
    %% 有说明文字，但目录中只显示标题（不含格式命令）
    %% 使用可选参数指定目录标题，显示标题后面加说明文字
    \@oldsubsection[#2]{#2 {\normalfont\kaishu\zihao{4}#1}}%
  }%
  \ifmnsfc@final\relax%
  \else\linenumbers\fi%% 恢复行号
  \songti\zihao{-4}%% 确保后续正文为宋体小四号
  \allowbreak
}

% 保留旧命令以保持向后兼容
\newcommand{\NSFCsection}[2][]{%
  \IfNoValueTF{#1}{%
    \section{#2}%
  }{%
    \section[#1]{#2}%
  }%
}

\newcommand{\NSFCsubsection}[2][]{%
  \IfNoValueTF{#1}{%
    \subsection{#2}%
  }{%
    \subsection[#1]{#2}%
  }%
}

%% 对于 amsrefs 包，修改 bibdiv 环境，避免创建 section
%% 在 AtBeginDocument 中修改，确保所有包都已加载
\AtBeginDocument{%
  %% 对于 amsrefs 包，重新定义 bibdiv 环境为空，避免创建 section
  %% 注意：amsrefs 的 MR 链接功能不依赖于 bibdiv 环境，
  %% 而是在处理 \bibliography 时直接添加到条目中，因此可以安全地重定义为空
  \@ifundefined{bibdiv}{}{%
    \renewenvironment{bibdiv}{}{}%
  }%
  %% 如果 amsrefs 未加载，给出提示
  \@ifundefined{bibdiv}{%
    \PackageWarning{mnsfc}{amsrefs package not loaded. %
      Please load it in your document.}%
  }{}%
}

% 参考文献命令
\newcommand{\mnsfcReferences}[1]{%
  \par\vspace{0.5\baselineskip}%
  \begin{center}%
    \kaishu\zihao{4}\textbf{参考文献}%
  \end{center}%
  \vspace{0.3\baselineskip}%
  \songti\zihao{-4}%
  \bibliography{#1}%
  \par\vspace{0.5\baselineskip}%
}

%% ==================================================
%% 字数统计功能（基于行号计算）
%% ==================================================
%% 创建计数器用于字数统计
\newcounter{mnsfc@totalpages}
\newcounter{mnsfc@totallines}
\newcounter{mnsfc@wordcount}
\newcounter{mnsfc@excludedlines}%% 被排除的短行计数

%% 标记短行（如 "无。"），使其不参与行号计数
\newcommand{\mnsfcExcludeLine}[1]{%
  \ifmnsfc@final%
    #1%
  \else%
    \stepcounter{mnsfc@excludedlines}%% 增加排除行计数
    \nolinenumbers%% 跳过这一行的行号显示
    #1\par%% 确保换行
    \linenumbers%% 恢复行号计数
    %% 注意：我们不在 here 减少行号计数器，而是在字数统计时减去排除的行数
    %% 这样可以保持行号显示的连续性
  \fi%
}

%% NSFC 格式：小四号，行距 22 磅
%% 每行约 35 个汉字（包括符号）
%% 计算公式：字数 = 行号 × 35字/行
\newcommand{\mnsfcWordCount}{%
  \ifmnsfc@final\relax%
  \else%
    %% 先获取行号值（在禁用行号之前）
    %% 获取总页数
    \setcounter{mnsfc@totalpages}{\getpagerefnumber{LastPage}}%
    %% 禁用行号（字数统计页面不需要行号）
    \nolinenumbers%
    %% 使用 \par 和 \vfill 代替 \clearpage，避免在文档末尾的 shipout 错误
    \par\vfill
    \begin{center}%
      \kaishu\zihao{4}\textbf{字数统计}%
    \end{center}%
    \vspace{0.5\baselineskip}%
    \songti\zihao{-4}%
    %% 尝试获取行号（如果 lineno 已启用）
    \@ifundefined{c@linenumber}{%
      %% 如果 lineno 未启用，使用页数估算
      \setcounter{mnsfc@totallines}%
        {\numexpr\value{mnsfc@totalpages}*35\relax}%
      \setcounter{mnsfc@wordcount}%
        {\numexpr\value{mnsfc@totallines}*35\relax}%
      \noindent 正文总页数：%
        \number\numexpr\value{mnsfc@totalpages}-1\relax 页\\
      \noindent 正文行数（估算）：约 \the\value{mnsfc@totallines} 行\\
      \noindent 正文字数（估算）：约 \the\value{mnsfc@wordcount} 字\\
      \vspace{0.3\baselineskip}%
      \noindent\textup{注：行号功能未启用，使用页数估算。}%
    }{%
      %% 如果 lineno 已启用，使用行号计算
      %% lineno 包使用 \c@linenumber 作为行号计数器
      %% 获取当前行号（在文档末尾调用时，这就是最后一行的行号）
      %% 减去被排除的短行数
      \setcounter{mnsfc@totallines}{\numexpr\value{linenumber}-1\relax}%
      \setcounter{mnsfc@wordcount}{\numexpr\value{mnsfc@totallines}*35\relax}%
      \noindent 正文总页数：%
        \number\numexpr\value{mnsfc@totalpages}-1\relax 页\\
      \noindent 正文行数（基于行号，已排除短行）：\the\value{mnsfc@totallines} 行\\
      \ifnum\value{mnsfc@excludedlines}>0%
        \noindent 已排除短行数：\the\value{mnsfc@excludedlines} 行\\
      \fi%
      \noindent 正文字数（估算）：约 \the\value{mnsfc@wordcount} 字\\
      \vspace{0.3\baselineskip}%
      \noindent\textup{注：字数 = 行数 × 35字/行。此统计为估算值，仅供参考。实际字数可能因排版、图表、公式等因素有所差异。}%
    }%
    \vfill
    \newpage%% 在内容之后分页，更安全
  \fi%
}
\usepackage{zhlipsum}
\ifmnsfc@final
\RenewDocumentCommand{\zhlipsum}{o}{}
\fi
%</package>
%    \end{macrocode}
%
% \section{示例文件}
%
% \subsection{使用示例文件}
%
%    \begin{macrocode}
%<*example-main>
%%!TeX encoding = utf8
\documentclass[UTF8,zihao=-4]{ctexart} %%for macOS set: fontset=macnew
%% ==模式切换 =====
\usepackage[final]{mnsfc} %%draft is default, final is for the final version
%% 加载 amsrefs 包，可以在这里控制选项
\usepackage[alphabetic,abbrev,lite,msc-links]{amsrefs}  %% 如果需要启用引文页码, 添加 backrefs 选项

\mnsfcProjectTitle{（填写项目名称）}
\mnsfcKeywordsZH{关键词1；关键词2；关键词3}
\mnsfcKeywordsEN{keywords1; keywords2; keywords3}
%% 选择研究属性（二选一）：
%% \mnsfcResearchTypeFree  %% 自由探索类基础研究
%% \mnsfcResearchTypeTarget %% 目标导向类基础研究
%% 或者使用自定义文本：
%% \mnsfcResearchType{自由探索类基础研究}
%% \mnsfcResearchType{目标导向类基础研究}
\mnsfcResearchTypeFree  %% 示例：选择自由探索类基础研究

 \begin{document}
   \mnsfcTitlePage
   \begin{mnsfcAbstractZH}
     这里是\emph{最终可提交版本}的中文摘要内容。
   \end{mnsfcAbstractZH}

   \begin{mnsfcAbstractEN}
     This is the English abstract for the NSFC proposal.
   \end{mnsfcAbstractEN}

  \mnsfcNote{正文开始...}
  \zhlipsum[1-2]
  %% 使用示例：subsubsection、paragraph、enumerate、itemize
   \subsubsection{三级标题示例}
   这是三级标题下的正文内容。三级标题使用编号格式 (1), (2), (3)...

   \paragraph{段落标题示例（不带编号）}
   这是段落标题下的正文内容。段落标题不带编号，仅作为强调使用。

   \subsubsection{列表环境示例}

   \textbf{有序列表（enumerate）示例：}
   \begin{enumerate}
     \item 第一层第一项：编号格式为 1), 2), 3)...
     \item 第一层第二项
     \begin{enumerate}
       \item 第二层第一项：编号格式为 a), b), c)...
       \item 第二层第二项
       \begin{enumerate}
         \item 第三层第一项：编号格式为 i), ii), iii)...
         \item 第三层第二项
       \end{enumerate}
     \end{enumerate}
     \item 第一层第三项
   \end{enumerate}

   \textbf{无序列表（itemize）示例：}
   \begin{itemize}
     \item 第一层第一项：使用实心圆点 •
     \item 第一层第二项
     \begin{itemize}
       \item 第二层第一项：使用空心圆点 ○
       \item 第二层第二项
       \begin{itemize}
         \item 第三层第一项：使用短横线 -
         \item 第三层第二项
       \end{itemize}
     \end{itemize}
     \item 第一层第三项
   \end{itemize}
   \subsubsection{参考文献}
   标准的参考文献数据可以从\href{https://mathscinet.ams.org/mrlookup}{美国数学会(MathSciNet)}获取。
   \cite{EinsteinBergmann1938AM}以及
   \cites{EinsteinBergmann1938AM,Hamilton1982JDG,Osserman1990AMM}
   \section[（为什么要开展此项研究，研究的科学技术价值如何）]{立项依据：}
   \zhlipsum[3-4]
   \begin{equation}\label{eq:NL}
     \int_a^b f'(x)\,dx=f(b)-f(a).
   \end{equation}
  从 \eqref{eq:NL}知,\ldots
  %%\bibliographystyle{plain}%
  \nocite{*}%
   \mnsfcReferences{mnsfc-refs}

   \section[（提纲不做限制，请按照研究工作的自身逻辑撰写。应提炼出特色与创新点、年度研究计划）]{研究内容：}
   \mnsfcNote{建议用"科学问题 → 国际现状 → 关键缺口 → 本项目目标"结构}
   \zhlipsum[5-6]
   \section{研究基础：}

   \subsection[（与本项目相关的研究工作积累和已取得的研究工作成绩，研究风险的应对措施等）；]{研究基础与可行性分析}
   \zhlipsum[7-8]

   \subsection[（包括已具备的实验条件，尚缺少的实验条件和拟解决的途径，包括利用国家实验室、全国重点实验室和部门重点实验室等研究基地的计划与落实情况）；]{工作条件}
   \zhlipsum[9-15]
   \subsection[（申请人和主要参与者正在承担的与本项目相关的科研项目情况，包括国家自然科学基金的项目和国家其他科技计划项目，要注明项目的资助机构、项目类别、批准号、项目名称、获资助金额、起止年月、与本项目的关系及负责的内容等）；]{正在承担的与本项目相关的科研项目情况}
  \mnsfcExcludeLine{无。}
  \subsection[（对申请人负责的前一个已资助期满的科学基金项目（项目名称及批准号）完成情况、后续研究进展及与本申请项目的关系加以详细说明。另附该项目的研究工作总结摘要（限500字）和相关成果详细目录）。]{完成国家自然科学基金项目情况}
  \mnsfcExcludeLine{无。}
  \section{其他需要说明的情况：}
  \subsection[申请人同年申请不同类型的国家自然科学基金项目情况（列明同年申请的其他项目的项目类型、项目名称信息，并说明与本项目之间的区别与联系；已收到自然科学基金委不予受理或不予资助决定的，无需列出）。]{}
  \mnsfcExcludeLine{无。}
  \subsection[具有高级专业技术职务（职称）的申请人或者主要参与者是否存在同年申请或者参与申请国家自然科学基金项目的单位不一致的情况；如存在上述情况，列明所涉及人员的姓名，申请或参与申请的其他项目的项目类型、项目名称、单位名称、上述人员在该项目中是申请人还是参与者，并说明单位不一致原因。]{}
  \mnsfcExcludeLine{无。}
  \subsection[具有高级专业技术职务（职称）的申请人或者主要参与者是否存在与正在承担的国家自然科学基金项目的单位不一致的情况；如存在上述情况，列明所涉及人员的姓名，正在承担项目的批准号、项目类型、项目名称、单位名称、起止年月，并说明单位不一致原因。]{}
  \mnsfcExcludeLine{无。}
  \subsection[申请人和主要参与者同年以不同专业技术职务（职称）申请或参与申请科学基金项目的情况（应详细说明原因）。]{}
  \mnsfcExcludeLine{无。}
  \subsection[其他。]{}
   \mnsfcNote{（包括但不限于使用以他人名义申报过的申请书；如有，请详细说明）。}
   \mnsfcNote{申请人在撰写本申请书时使用生成式人工智能的情况，请详细说明申请书中使用的位置和内容。}
   \mnsfcNote{申请人在撰写申请书时，如果借助生成式人工智能技术跟踪研究动态、收集整理参考文献，必须由人工核实生成式人工智能生成信息和参考文献的真实性和准确性。申请人应当对使用的生成内容负责，应当全面如实声明使用情况，按照《人工智能生成合成内容标识办法》《负责任研究行为规范指引（2023）》等国家有关规定对相关内容进行标识，标识的方式包括但不限于在文本的起始和末尾适当位置添加相应的文字提示等。不得使用由生成式人工智能直接生成的申请书，不得使用未经核实的生成内容。}
  \mnsfcExcludeLine{无。}
  %% 字数统计（仅在 draft 模式下显示）
  \mnsfcWordCount
\end{document}
%</example-main>
%    \end{macrocode}
%
% \subsection{示例参考文献数据}
% 标准的参考文献数据可以通过\href{https://mathscinet.ams.org/mrlookup}{美国数学会(MathSciNet)}获取。
%    \begin{macrocode}
%<*example-bib>
%% 参考文献示例文件
%% 使用标准 BibTeX 格式
@article {Hamilton1982JDG,
    AUTHOR = {Hamilton, Richard S.},
     TITLE = {Three-manifolds with positive {R}icci curvature},
   JOURNAL = {J. Differential Geometry},
  FJOURNAL = {Journal of Differential Geometry},
    VOLUME = {17},
      YEAR = {1982},
    NUMBER = {2},
     PAGES = {255--306},
      ISSN = {0022-040X,1945-743X},
   MRCLASS = {53C25 (35K55 58G30)},
  MRNUMBER = {664497},
MRREVIEWER = {J.\ L.\ Kazdan},
       URL = {http://projecteuclid.org/euclid.jdg/1214436922},
}
@article {Osserman1990AMM,
    AUTHOR = {Osserman, Robert},
     TITLE = {Curvature in the eighties},
   JOURNAL = {Amer. Math. Monthly},
  FJOURNAL = {American Mathematical Monthly},
    VOLUME = {97},
      YEAR = {1990},
    NUMBER = {8},
     PAGES = {731--756},
      ISSN = {0002-9890,1930-0972},
   MRCLASS = {53-01 (53-02 53C42)},
  MRNUMBER = {1072814},
MRREVIEWER = {M.\ do Carmo},
       DOI = {10.2307/2324577},
       URL = {https://doi.org/10.2307/2324577},
}
@article {EinsteinBergmann1938AM,
    AUTHOR = {Einstein, A. and Bergmann, P.},
     TITLE = {On a generalization of {K}aluza's theory of electricity},
   JOURNAL = {Ann. of Math. (2)},
  FJOURNAL = {Annals of Mathematics. Second Series},
    VOLUME = {39},
      YEAR = {1938},
    NUMBER = {3},
     PAGES = {683--701},
      ISSN = {0003-486X,1939-8980},
   MRCLASS = {DML},
  MRNUMBER = {1503432},
       DOI = {10.2307/1968642},
       URL = {https://doi.org/10.2307/1968642},
}
%</example-bib>
%    \end{macrocode}
%
%\CheckSum{0}
% \Finale
\endinput
